\section{Prime Numbers}

\subsection{Primes}

\begin{defbox}
   An integer $p \in (\ZZ \setminus \{0, \pm 1\}$ is said to be \textbf{prime} if its only divisors are $1$ and $p$. A \textbf{prime number} is a natural number which is prime, i.e. a number $p \geq 2$ such that its divisors are $1$ and $p$. 
\end{defbox}

Note that by definition, $1$ is not a prime number. 

\subsection{The Sieve of Eratosthenes}

To find prime numbers, we can use the \textit{Sieve of Eratosthenes}, which is a slow but reliable algorithm to find all prime numbers from $2$ to a natural number $N$. 

Suppose, for sake of exposition, we want to find all the prime numbers between $2$ and $50$. We first write out all numbers between $2$ and $50$:
\[
\begin{array}{rrrrrrrrr}
2 & 3 & 4 & 5 & 6 & 7 & 8 & 9 & 10 \\
11 & 12 & 13 & 14 & 15 & 16 & 17 & 18 & 19 \\
20 & 21 & 22 & 23 & 24 & 25 & 26 & 27 & 28 \\
29 & 30 & 31 & 32 & 33 & 34 & 35 & 36 & 37 \\
38 & 39 & 40 & 41 & 42 & 43 & 44 & 45 & 46 \\
47 & 48 & 49 & 50
\end{array}
\]
We start by highlighting the smallest number--in this case it is $2$. $2$ is a prime number. We then cross out all multiples of $2$, as they are not prime:
\[
\begin{array}{rrrrrrrrr}
\boxed{2} & 3 & \cancel{4} & 5 & \cancel{6} & 7 & \cancel{8} &
\cancel{9} & \cancel{10} \\
11 & \cancel{12} & 13 & \cancel{14} & \cancel{15} & \cancel{16} &
17 & \cancel{18} & 19 \\
\cancel{20} & \cancel{21} & \cancel{22} & 23 & \cancel{24} &
\cancel{25} & \cancel{26} & \cancel{27} & \cancel{28} \\
29 & \cancel{30} & 31 & \cancel{32} & \cancel{33} & \cancel{34} &
\cancel{35} & \cancel{36} & 37 \\
\cancel{38} & \cancel{39} & \cancel{40} & 41 & \cancel{42} &
43 & \cancel{44} & \cancel{45} & \cancel{46} \\
47 & \cancel{48} & \cancel{49} & \cancel{50}
\end{array}
\]
Now, we move to highlight the next smallest number: $3$. We now cross out all multiples of $3$ (we can skip numbers that have already been crossed out):
\[
\begin{array}{rrrrrrrrr}
\boxed{2} & \boxed{3} & \cancel{4} & 5 & \cancel{6} & 7 &
\cancel{8} & \cancel{9} & \cancel{10} \\
11 & \cancel{12} & 13 & \cancel{14} & \cancel{15} & \cancel{16} &
17 & \cancel{18} & 19 \\
\cancel{20} & \cancel{21} & \cancel{22} & 23 & \cancel{24} &
25 & \cancel{26} & \cancel{27} & \cancel{28} \\
29 & \cancel{30} & 31 & \cancel{32} & \cancel{33} & \cancel{34} &
35 & \cancel{36} & 37 \\
\cancel{38} & \cancel{39} & \cancel{40} & 41 & \cancel{42} &
43 & \cancel{44} & \cancel{45} & \cancel{46} \\
47 & \cancel{48} & \cancel{49} & \cancel{50}
\end{array}
\]
The next smallest number, $4$, has been crossed out (hence is not prime), so we move to highlight $5$, which is prime. We now cross out all multiples of $5$:
\[
\begin{array}{rrrrrrrrr}
\boxed{2} & \boxed{3} & \cancel{4} & \boxed{5} & \cancel{6} & 7 &
\cancel{8} & \cancel{9} & \cancel{10} \\
11 & \cancel{12} & 13 & \cancel{14} & \cancel{15} & \cancel{16} &
17 & \cancel{18} & 19 \\
\cancel{20} & \cancel{21} & \cancel{22} & 23 & \cancel{24} &
\cancel{25} & \cancel{26} & \cancel{27} & \cancel{28} \\
29 & \cancel{30} & 31 & \cancel{32} & \cancel{33} & \cancel{34} &
\cancel{35} & \cancel{36} & 37 \\
\cancel{38} & \cancel{39} & \cancel{40} & 41 & \cancel{42} &
43 & \cancel{44} & \cancel{45} & \cancel{46} \\
47 & \cancel{48} & \cancel{49} & \cancel{50}
\end{array}
\]
Continuing this algorithm until all numbers have been either crossed out or highlighed, we get the prime numbers up to $50$:
\[
\boxed{2}\quad \boxed{3}\quad \boxed{5}\quad \boxed{7}\quad
\boxed{11}\quad \boxed{13}\quad \boxed{17}\quad \boxed{19}\quad
\boxed{23}\quad \boxed{29}\quad \boxed{31}\quad \boxed{37}\quad
\boxed{41}\quad \boxed{43}\quad \boxed{47}
\]

\subsection{Properties of Prime Numbers}

\begin{thmbox}[Characterisation of Primes]
   Let $p \geq 2$ be a natural number. $p$ is prime if and only if for all $a$, $b \in \ZZ$, if $p \mid ab$, then $p \mid a$ or $p \mid b$. 
\end{thmbox}
\begin{proof}
   Assume that $p$ is prime and that for all $a$, $b \in \ZZ$, $p \mid ab$. Assume $p \nmid a$, then $\gcd(a,p) = 1$ since $p$ is prime. Thus, by Corollary 8.1.3, $p \mid b$. Conversely, suppose that for all $a$, $b \in \ZZ$, $p \mid ab$ implies $p \mid a$ and $p \mid b$. We suppose $p = st$ for some $s$, $t \in \NN^+$ and show that $s = 1$ or $t = 1$. Since $p \mid st$, $p \mid s$, or $p \mid t$, say $p \mid s$. Then $s = ps'$ hence $p = st = s't$ so $s't = 1$, hence $s' = t = 1$. 
\end{proof}

\begin{corrbox}
   If $p \geq 2$ is prime and $p \mid a_1a_2\dots a_n$ for $a_1, \dots, a_n \in \ZZ$, then there exists an $i \leq n$ such that $p \mid ai$.
\end{corrbox}
\begin{proof}
   Can be done trivially by induction of $n$ using the Characterisation of Primes. 
\end{proof}

\begin{thmbox}
   Every natural number $n \geq 1$ is a product of primes.
\end{thmbox}
\begin{proof}
   Shown in homework.
\end{proof}

\begin{corrbox}
   Every natural number $n \geq 2$ is divisible by a prime.
\end{corrbox}
\begin{proof}
   TBD.
\end{proof}

\begin{corrbox}[Infinitude of Primes]
   There are infinitely many prime numbers.
\end{corrbox}
\begin{proof}
   Suppose by contradictoion that there are finitely many primes $p_1$, $\dots$, $p_k$ for some $k \geq 1$. Let $n := p_1p_2\dots p_k + 1$. Then there is a prime $q \mid n$ by the previous corrolary, so $q = p_i$ for some $i \in \{1, \dots, k\}$ since $p_1, \dots, p_k$ are all the primes. But then $1 = n - p_1\dots p_k$ is divisible by $q = p_i$ since $q \mid a$ and $q \mid p_1 \dots p_k$, which is a contradiction.
\end{proof}

\begin{thmbox}[Fundamental Theorem of Arithmetic]
   Every nonzero integer $n$ is a unique product of prime integers, i.e. $n = \sigma p_1\dots p_k$, where $\sigma = \pm 1$ and $2 \leq q_1 \leq \dots \leq q_k$ are primes.
\end{thmbox}
\begin{proof}
   The existence of such a representation follows from Theorem 9.3.4. We only need to show the uniqueness of such a representation.

   Let $\sigma p_1 \dots p_k = n = \sigma' q_1 \dots q_m$ where $p_i$ and $p_j$ are prime numbers ordered $2 \leq p_1 \leq \dots \leq p_k$ and $2 \leq q_1 \leq \dots \leq q_m$. The sign must be the same, i.e. $\sigma = \sigma'$, so assume without loss of generality that $n \geq 0$ and $p_1 \dots p_k = n = q_1 \dots q_m$. We prove that $k = m$ and $p_i = q_i$ for all $1 \leq i \leq k$. We can show this holds by strong induction on $n \geq 1$.

   If $n = 1$, then $k = 0 = m$ and we are done. If we suppose $n \geq 2$, then $k$, $m \geq$ and without loss of generality suppose $p_1 \leq q_1$. Then $p_1 \mid q_1 \dots q_m$ so $p_i \mid q_j$ for some $1 \leq j \leq m$. But $q_j$ is prime so $p_1 = q_j$. Since $q_1 \leq \dots q_j$ and $p_1 \leq q_1$, we get $p_1 = q_1 = \dots = q_j$, so $n' := p_2p_3\dots p_k = q_2q_3 \dots q_m$. By the induction hypothesis on $n'$, we get $k = m$ and $p_i = q_i$ for all $i \in [2, k]$, hence $k = m$ and $p_i = q_i$ for all $i \in [1,k]$, whence the representation is unique.
\end{proof}

\begin{propbox}
   Any non-zero rational number $q$ can be written as
   \[
   q = \sigma p_1^{a_1} \dots p_m^{a_m},
   \]
   where $\sigma = \pm 1$, the $p_i$ are distinct prime numbers, and
   $a_1, \dots, a_m$ are non-zero integers, possibly negative. Moreover,
   this expression is unique up to reordering the primes.
\end{propbox}

\begin{proof}
   Since $q$ is rational, there exists $a$, $b \in \ZZ$ with $b \neq 0$ such that
   \[
   q = \frac{a}{b}.
   \]
   By the Fundamental Theorem of Arithmetic, we can write
   \[
   a = \sigma_a r_1^{s_1} \dots r_n^{s_n},
   \qquad
   b = \sigma_b r_1^{t_1} \dots r_n^{t_n},
   \]
   where $\sigma_a, \sigma_b \in \{-1,1\}$, the $r_i$ are distinct prime
   numbers, and $s_i, t_i$ are non-negative integers. We allow some of the
   exponents to be zero so that the same list of primes can be used in both
   factorizations. Therefore,
   \[
   q = \frac{\sigma_a}{\sigma_b}
   r_1^{s_1-t_1}\dots r_n^{s_n-t_n}.
   \]
   Omitting the primes for which $s_i-t_i=0$, and calling the remaining
   primes $p_1,\dots,p_m$, we obtain
   \[
   q = \sigma p_1^{a_1}\dots p_m^{a_m},
   \]
   where $\sigma=\sigma_a/\sigma_b \in \{-1,1\}$ and each $a_i$ is a
   non-zero integer.

   Now suppose that we have two such expressions for $q$. By allowing zero
   exponents, it is enough to consider
   \[
   q = \sigma p_1^{a_1}\dots p_m^{a_m}
   = \sigma' p_1^{a'_1}\dots p_m^{a'_m},
   \]
   where the $p_i$ are distinct prime numbers. Since $\sigma$ and $\sigma'$
   determine the sign of $q$, we must have $\sigma=\sigma'$. Dividing through
   by the second expression gives
   \[
   1 = p_1^{c_1}\dots p_m^{c_m},
   \]
   where $c_i=a_i-a'_i$. We need to show that $c_i=0$ for all $i$.

   By reordering the primes, we may assume that $c_1,\dots,c_t$ are
   negative and that $c_{t+1},\dots,c_m$ are non-negative. Hence,
   \[
   p_1^{-c_1}\dots p_t^{-c_t}
   =
   p_{t+1}^{c_{t+1}}\dots p_m^{c_m}.
   \]
   Both sides are now products of primes with non-negative exponents. By
   unique factorization for integers, and since all the primes are distinct,
   every exponent must be zero. Thus $c_i=0$ for all $i$, so
   \[
   a_i=a'_i
   \]
   for all $i$. Therefore, the expression is unique up to reordering the
   primes.
\end{proof}

